\documentclass[12pt,a4paper]{article}
\usepackage[utf8]{inputenc}
\usepackage[spanish]{babel}
\usepackage{geometry}
\usepackage{xcolor}
\usepackage{listings}
\usepackage{hyperref}
\usepackage{graphicx}
\usepackage{float}

\geometry{top=2.5cm, bottom=2.5cm, left=3cm, right=3cm}

\definecolor{codegreen}{rgb}{0,0.6,0}
\definecolor{codegray}{rgb}{0.5,0.5,0.5}
\definecolor{codepurple}{rgb}{0.58,0,0.82}
\definecolor{backcolour}{rgb}{0.95,0.95,0.92}

\lstdefinestyle{sqlstyle}{
    backgroundcolor=\color{backcolour},   
    commentstyle=\color{codegreen},
    keywordstyle=\color{blue},
    numberstyle=\tiny\color{codegray},
    stringstyle=\color{codepurple},
    basicstyle=\ttfamily\footnotesize,
    breakatwhitespace=false,         
    breaklines=true,                 
    captionpos=b,                    
    keepspaces=true,                 
    numbers=left,                    
    numbersep=5pt,                  
    showspaces=false,                
    showstringspaces=false,
    showtabs=false,                  
    tabsize=2,
    language=SQL
}
\lstset{style=sqlstyle}

\title{\textbf{Documentación Técnica y Arquitectura de Software: Steam BDA} \\ \Large{Implementación de un E-Commerce con Motor Transaccional en Oracle Database}}
\author{Proyecto Final de Bases de Datos Avanzadas}
\date{\today}

\begin{document}

\maketitle
\tableofcontents
\newpage

\section{Introducción}
El presente documento expone la arquitectura técnica y el diseño transaccional del proyecto "Steam BDA", una plataforma web de distribución digital de videojuegos. Si bien el proyecto cuenta con una interfaz gráfica moderna desarrollada en Python (Flask) y HTML/CSS, el objetivo de esta documentación es demostrar que el verdadero procesamiento (Backend), las reglas de negocio, la integridad de los datos y el rendimiento están implementados enteramente en el motor de \textbf{Oracle Database}.

Flask funciona únicamente como un intermediario (View/Controller), mientras que Oracle actúa como el cerebro del sistema aplicando conceptos avanzados como \textit{Triggers, Stored Procedures, List Partitioning y Bitmap Indexes}.

\section{Perfiles de Usuario y sus Módulos}
El sistema implementa un modelo de control de acceso basado en roles (RBAC) gestionado directamente desde la base de datos. Existen tres roles principales:

\subsection{1. Administrador (Rol: ADMIN)}
Es el supervisor de la plataforma. Su módulo no le permite comprar juegos, sino administrar la integridad del ecosistema comercial.
\begin{itemize}
    \item \textbf{Módulo de Aprobaciones (Solicitudes):} Visualiza los juegos en estado pendiente (\texttt{PE}) subidos por las empresas. Puede aprobarlos para que sean públicos o denegarlos.
    \item \textbf{Módulo de Revocación/Restauración:} Puede realizar un \textit{Soft Delete} (quitar del catálogo) a juegos ya aprobados si infringen normas, utilizando el procedimiento \texttt{sp\_revocar\_juego}. Este proceso ahora elimina en cascada el juego de los carritos activos y \textbf{revoca el acceso} de las bibliotecas de los usuarios, respaldando dicha información en la tabla \texttt{biblioteca\_revocada}. Luego puede usar el botón "Restaurar" mediante \texttt{sp\_restaurar\_juego}, el cual devuelve el juego al catálogo y re-establece automáticamente las licencias a los usuarios afectados.
    \item \textbf{Módulo de Eliminación Permanente:} Si un juego revocado necesita ser borrado físicamente, el administrador puede ejecutar \texttt{sp\_eliminar\_juego\_permanente}, el cual limpia en cascada todas las dependencias (calificaciones, bibliotecas revocadas, detalles de pedido, etc.) y realiza un borrado físico final, dejando un rastro en auditoría.
    \item \textbf{Módulo de Auditoría:} Consulta la tabla \texttt{auditoria} que es alimentada automáticamente por los \textit{Triggers} del sistema sin intervención humana.
    \item \textbf{Módulo de Actividad de Usuarios:} Visualiza reportes gerenciales a través de la vista analítica \texttt{v\_actividad\_compras}.
\end{itemize}

% PARA AGREGAR IMAGENES: Descomenta las siguientes líneas y asegúrate de tener la imagen en la misma carpeta.
% \begin{figure}[H]
% \centering
% \includegraphics[width=0.9\textwidth]{tu_captura_admin.png} 
% \caption{Módulo Administrativo y Gestión de Solicitudes}
% \end{figure}

\subsection{2. Desarrollador / Empresa (Rol: EMPRESA)}
Representa a los creadores de contenido (Estudios de videojuegos) que buscan publicar sus obras.
\begin{itemize}
    \item \textbf{Módulo de Subida de Juegos:} Interfaz que permite a la empresa publicar nuevos títulos indicando título, precio, descripción, y subir archivos multimedia (Carátulas y Trailers de video) eligiendo entre almacenamiento puro en BD (BLOB) o en el disco del servidor (RUTA).
    \item \textbf{Módulo "Mis Juegos":} Panel administrativo para la empresa donde puede listar todos sus títulos, visualizar su estado de aprobación y modificar sus datos (precio, descripción, multimedia). Esta edición está vigilada por \textit{Triggers} de Oracle para proteger a los consumidores. 
    \item \textbf{Módulo de Relación Muchos a Muchos (Géneros):} La interfaz presenta una lista múltiple de géneros. Flask captura estos datos y mediante una instrucción \texttt{RETURNING}, Oracle inserta dinámicamente los registros en la tabla puente \texttt{videojuego\_genero}. Todo juego nuevo entra forzosamente en estado Pendiente (\texttt{PE}).
\end{itemize}

\subsection{3. Usuario Final (Rol: USUARIO)}
Es el consumidor y cliente final de la plataforma web.
\begin{itemize}
    \item \textbf{Módulo del Catálogo Público:} Una galería visual amplia (alimentada por la vista \texttt{v\_catalogo\_publico}) donde solo visualiza los juegos con estado Aprobado (\texttt{AP}).
    \item \textbf{Módulo del Carrito y Pagos (Checkout):} Carrito de compras temporal en la base de datos que, al confirmarse, dispara el procedimiento de transacción ACID pesada (\texttt{sp\_procesar\_compra}).
    \item \textbf{Módulo de Biblioteca Virtual:} Espacio personal tipo estantería donde residen permanentemente las licencias adquiridas.
    \item \textbf{Módulo de Calificaciones (Reseñas):} Componente visual incrustado en la Biblioteca que permite otorgar de 1 a 5 estrellas a un juego. Estos votos se almacenan mediante \textit{List Partitioning} físico.
\end{itemize}

% \begin{figure}[H]
% \centering
% \includegraphics[width=0.9\textwidth]{tu_captura_catalogo.png} 
% \caption{Catálogo Público consumiendo Vistas de Oracle}
% \end{figure}

\section{Arquitectura de la Aplicación (Frontend y Middleware)}
El sistema web se construyó utilizando el framework Flask (Python) implementando un patrón de diseño basado en \textbf{MVC} (Modelo-Vista-Controlador), pero con el "Modelo" pesado delegado a la base de datos Oracle. Esta estructura modular (utilizando \textit{Blueprints} de Flask) garantiza que el código sea mantenible, escalable y ordenado.

\subsection{Estructura de Archivos Principales}
\begin{itemize}
    \item \textbf{\texttt{app.py}:} Es el punto de entrada principal (Entry Point). Su responsabilidad es inicializar la aplicación, registrar los "Blueprints" (módulos), configurar las sesiones de los usuarios y levantar el servidor web local.
    \item \textbf{\texttt{database.py}:} Contiene la configuración centralizada y el Pool de Conexiones usando el módulo nativo \texttt{oracledb} de Python. Todas las consultas viajan por este archivo.
\end{itemize}

\subsection{Controladores (Carpeta \texttt{controllers/})}
Contienen la lógica de enrutamiento (Middleware) y actúan como intermediarios entre las acciones del usuario (clics en botones) y los Procedimientos de Oracle.
\begin{itemize}
    \item \textbf{\texttt{auth\_controller.py}:} Gestiona las sesiones seguras, el Inicio de Sesión y el Registro de nuevos usuarios o empresas.
    \item \textbf{\texttt{admin\_controller.py}:} Protegido por validación de roles. Llama a las vistas de auditoría, las funciones de aprobar, denegar, revocar o restaurar juegos.
    \item \textbf{\texttt{games\_controller.py}:} Maneja la carga pesada. Captura la subida de nuevos títulos (procesando imágenes y videos BLOB), extrae los géneros relacionales y se encarga de mostrar el Catálogo Público leyendo la vista \texttt{v\_catalogo\_publico}.
    \item \textbf{\texttt{cart\_controller.py}:} Administra el estado del carrito en la base de datos y detona el procedimiento transaccional \texttt{sp\_procesar\_compra} para efectuar el pago y poblar la biblioteca del usuario.
    \item \textbf{\texttt{reviews\_controller.py}:} Captura las calificaciones enviadas desde la biblioteca y las inserta en la tabla particionada de Oracle.
\end{itemize}

\subsection{Vistas y Plantillas (Carpeta \texttt{templates/})}
Archivos en formato HTML/Jinja2 que construyen la interfaz visual usando la técnica \textit{Glassmorphism} e inyectan dinámicamente los datos traídos desde Oracle.
\begin{itemize}
    \item \textbf{\texttt{base.html}:} La plantilla maestra. Contiene la barra de navegación inteligente que muestra enlaces distintos dependiendo del rol activo en la sesión (Usuario, Empresa o Admin).
    \item \textbf{\texttt{catalogo.html}:} Una cuadrícula (\textit{Grid}) inmersiva que despliega el arte, los géneros y el promedio de estrellas de cada juego aprobado.
    \item \textbf{\texttt{subir\_juego.html}:} Formulario para que las Empresas publiquen nuevos títulos, adjuntando contenido multimedia (BLOB) y géneros relacionales.
    \item \textbf{\texttt{biblioteca.html}:} Diseñada como una estantería virtual. Solo muestra las portadas de los juegos comprados; al pasar el ratón (hover) revela un formulario minimalista para dejar calificaciones y reseñas.
    \item \textbf{\texttt{admin/pendientes.html}:} Panel de control de 2 vías del Administrador, donde puede Aprobar/Denegar solicitudes (estado \texttt{PE}) o Restaurar juegos castigados (estado \texttt{RE}).
    \item \textbf{\texttt{admin/auditoria.html} y \texttt{actividad.html}:} Tablas de datos en bruto sacadas de las Vistas de Oracle para control gerencial.
\end{itemize}

\subsection{Estilos Globales (Carpeta \texttt{static/})}
\begin{itemize}
    \item \textbf{\texttt{style.css}:} Hoja de estilos en cascada única que le otorga a la plataforma su apariencia \textit{Premium} (temática "Cool-Tone", fondos de cristal desenfocado, colores cianes/azules intensos, botones con bordes luminosos y cuadrículas adaptables). Evita el uso de librerías externas genéricas para tener control de píxel.
\end{itemize}

\section{Diseño Físico y Almacenamiento Avanzado}
Para garantizar la eficiencia a largo plazo y la prevención de fragmentación en el disco, se emplearon técnicas propias de un Administrador de Bases de Datos (DBA).

\subsection{Tablespaces y Parámetros PCTFREE / PCTUSED}
Se separó la carga de Entrada/Salida creando \texttt{ts\_datos} para datos puros y \texttt{ts\_indices} para índices. Además, se personalizaron los bloques físicos en las sentencias de creación:

\begin{lstlisting}
-- Ejemplo de la tabla Videojuego, reservando 20% del bloque para 
-- futuras actualizaciones de estado y previniendo "Row Migration".
CREATE TABLE videojuego (
    id_juego NUMBER GENERATED BY DEFAULT AS IDENTITY PRIMARY KEY,
    titulo VARCHAR2(200) NOT NULL,
    precio NUMBER(10,2) CHECK (precio >= 0),
    estado VARCHAR2(2) DEFAULT 'PE' CHECK (estado IN ('PE','AP','RE'))
    -- Resto de columnas omitidas por brevedad...
) TABLESPACE ts_datos PCTFREE 20;

-- Ejemplo de Auditoría, minimizando el espacio libre a solo el 5% 
-- debido a su naturaleza de "Solo Inserción" (Append-only).
CREATE TABLE auditoria (
    accion VARCHAR2(100),
    fecha DATE DEFAULT SYSDATE
) TABLESPACE ts_datos PCTFREE 5;
\end{lstlisting}

\subsection{Particionamiento Físico (List Partitioning)}
Para agilizar los reportes sobre los juegos mejores o peores valorados, la tabla \texttt{calificaciones\_part} se seccionó físicamente usando particionamiento de listas basado en la puntuación del 1 al 5:

\begin{lstlisting}
CREATE TABLE calificaciones_part (
    id_usuario NUMBER,
    id_juego NUMBER,
    puntuacion NUMBER(1) CHECK (puntuacion BETWEEN 1 AND 5),
    UNIQUE(id_usuario, id_juego)
) TABLESPACE ts_datos PCTFREE 10 PCTUSED 40
PARTITION BY LIST (puntuacion) (
    PARTITION p1 VALUES (1) TABLESPACE ts_datos,
    PARTITION p2 VALUES (2) TABLESPACE ts_datos,
    PARTITION p3 VALUES (3) TABLESPACE ts_datos,
    PARTITION p4 VALUES (4) TABLESPACE ts_datos,
    PARTITION p5 VALUES (5) TABLESPACE ts_datos
);
\end{lstlisting}

\subsection{Índices Especializados (Bitmap)}
Se implementó un índice en formato de mapa de bits (\texttt{BITMAP INDEX}) sobre la columna \texttt{estado} de la tabla \texttt{videojuego}. A diferencia de los B-Trees convencionales, los Bitmap son ideales para acelerar consultas sobre columnas de muy baja cardinalidad (donde solo hay valores PE, AP, RE).
\begin{lstlisting}
CREATE BITMAP INDEX idx_videojuego_estado_bitmap 
ON videojuego(estado) TABLESPACE ts_indices;
\end{lstlisting}

\section{Implementación Transaccional del CRUD (PL/SQL)}
En lugar de procesar la lógica frágil en Python, las transacciones críticas y de modificación múltiple (CRUD) se realizan exclusivamente mediante Procedimientos Almacenados, garantizando las propiedades ACID.

\subsection{Procedimiento de Compra (Transacción Multitabla)}
El proceso \texttt{sp\_procesar\_compra} transforma el estado del carrito en un pedido final, moviendo licencias y calculando montos en una única transacción controlada.

\begin{lstlisting}
CREATE OR REPLACE PROCEDURE sp_procesar_compra (
    p_usuario_id IN NUMBER,
    p_resultado OUT VARCHAR2
) AS
    v_total NUMBER := 0;
    v_pedido_id NUMBER;
BEGIN
    SELECT NVL(SUM(v.precio), 0) INTO v_total
    FROM carrito c JOIN videojuego v ON c.id_juego = v.id_juego
    WHERE c.id_usuario = p_usuario_id;
    
    IF v_total = 0 THEN
        p_resultado := 'ERROR: Carrito vacio.';
        RETURN;
    END IF;

    -- 1. Insertar la cabecera del pedido
    INSERT INTO pedido (id_usuario, total, estado)
    VALUES (p_usuario_id, v_total, 'COMPLETADO')
    RETURNING id_pedido INTO v_pedido_id;

    -- 2. Recorrer el carrito e insertar en Detalle y Biblioteca
    FOR reg IN (SELECT c.id_juego, v.precio FROM carrito c 
                JOIN videojuego v ON c.id_juego = v.id_juego 
                WHERE c.id_usuario = p_usuario_id) LOOP
                
        INSERT INTO detalle_pedido (id_pedido, id_juego, precio_pagado)
        VALUES (v_pedido_id, reg.id_juego, reg.precio);
        
        INSERT INTO biblioteca (id_usuario, id_juego)
        VALUES (p_usuario_id, reg.id_juego);
    END LOOP;

    -- 3. Limpiar carrito y consolidar (ACID)
    DELETE FROM carrito WHERE id_usuario = p_usuario_id;
    COMMIT;
    p_resultado := 'OK';
EXCEPTION
    WHEN OTHERS THEN
        ROLLBACK;
        RAISE;
END;
\end{lstlisting}

\subsection{Revocación de Juegos (Borrado Lógico / Soft Delete)}
Para respetar las leyes de derechos digitales de los consumidores, en lugar de ejecutar un \texttt{DELETE} destructivo, se programó un procedimiento de borrado lógico que oculta el juego del catálogo público, pero mantiene el historial financiero y de la biblioteca intacto.

\begin{lstlisting}
CREATE OR REPLACE PROCEDURE sp_revocar_juego (
    p_juego_id  IN NUMBER,
    p_admin_id  IN NUMBER,
    p_resultado OUT VARCHAR2
) AS
    v_estado VARCHAR2(2);
BEGIN
    SELECT estado INTO v_estado FROM videojuego WHERE id_juego = p_juego_id;
    
    IF v_estado != 'AP' THEN
        p_resultado := 'ERROR: El juego no esta publicado.';
        RETURN;
    END IF;
    
    -- Respaldar las bibliotecas afectadas antes de borrarlas
    INSERT INTO biblioteca_revocada (id_biblioteca, id_usuario, id_juego, fecha_adquisicion)
    SELECT id_biblioteca, id_usuario, id_juego, fecha_adquisicion
    FROM biblioteca WHERE id_juego = p_juego_id;
    
    -- Eliminar de bibliotecas y carritos
    DELETE FROM biblioteca WHERE id_juego = p_juego_id;
    DELETE FROM carrito WHERE id_juego = p_juego_id;
    
    UPDATE videojuego 
    SET estado = 'RE', aprobado_por = p_admin_id, fecha_aprobacion = SYSDATE
    WHERE id_juego = p_juego_id;
    
    COMMIT;
    p_resultado := 'OK';
END;
\end{lstlisting}

\section{Reglas de Negocio Estrictas (Triggers)}
Los Disparadores (Triggers) operan de manera subterránea para asegurar que ningún actor viole las reglas de negocio empresariales, lanzando errores que luego son capturados limpiamente por el frontend de Flask.

\subsection{Validación Lógica de Reseñas}
La regla dicta que "nadie puede reseñar un producto que no haya comprado". Este Trigger se ejecuta en el pre-proceso de inserción, bloqueando la base de datos si la condición es falsa.
\begin{lstlisting}
CREATE OR REPLACE TRIGGER trg_validar_compra_resena
BEFORE INSERT ON calificaciones_part
FOR EACH ROW
DECLARE
    v_tiene_juego NUMBER;
BEGIN
    SELECT COUNT(*) INTO v_tiene_juego FROM biblioteca 
    WHERE id_usuario = :NEW.id_usuario AND id_juego = :NEW.id_juego;
    
    IF v_tiene_juego = 0 THEN
        RAISE_APPLICATION_ERROR(-20002, 'ERROR: Solo puedes calificar juegos que has comprado y tienes en tu biblioteca.');
    END IF;
END;
\end{lstlisting}

\subsection{Protección al Consumidor (Inmutabilidad)}
Para evitar fraudes post-venta (ej. un desarrollador cambiando el nombre de un juego exitoso por publicidad de terceros), este Trigger impide modificaciones estructurales una vez que ha sido aprobado (\texttt{estado = 'AP'}).
\begin{lstlisting}
CREATE OR REPLACE TRIGGER trg_proteger_juego_aprobado
BEFORE UPDATE ON videojuego
FOR EACH ROW
BEGIN
    IF :OLD.estado = 'AP' THEN
        IF :NEW.titulo != :OLD.titulo THEN
            RAISE_APPLICATION_ERROR(-20004, 'ERROR: No se puede cambiar el titulo de un juego que ya esta publicado y aprobado.');
        END IF;
    END IF;
END;
\end{lstlisting}

\section{Vistas Complejas y Reporteo}
Para que la interfaz web reciba la información pre-digerida, se diseñaron Vistas SQL que empaquetan sub-consultas, conteos y funciones.

\subsection{Vista de Catálogo (Uso de Funciones y LISTAGG)}
Para abstraer la relación de "muchos a muchos" de la tabla de géneros, se utilizó la instrucción avanzada \texttt{LISTAGG}, la cual pivotea las filas relacionales integrándolas en una única cadena de texto separada por comas, eliminando la necesidad de iteraciones costosas en Python.

\begin{lstlisting}
CREATE OR REPLACE VIEW v_catalogo_publico AS
SELECT 
    v.id_juego,
    v.titulo,
    v.precio,
    u.username AS empresa,
    -- Invocando a una Funcion almacenada
    fn_promedio_calificacion(v.id_juego) AS calificacion_promedio,
    -- Subconsulta Escalar para totales
    (SELECT COUNT(*) FROM calificaciones_part cp WHERE cp.id_juego = v.id_juego) AS num_resenas,
    -- LISTAGG para consolidacion de Generos
    (SELECT LISTAGG(g.nombre, ', ') WITHIN GROUP (ORDER BY g.nombre)
     FROM videojuego_genero vg
     JOIN genero g ON vg.id_genero = g.id_genero
     WHERE vg.id_juego = v.id_juego) AS generos
FROM videojuego v
JOIN usuario u ON v.id_empresa = u.id_usuario
WHERE v.estado = 'AP';
\end{lstlisting}

\section{Conclusión}
La arquitectura tecnológica de "Steam BDA" evidencia la madurez de \textbf{Oracle Database} como pilar para ecosistemas de gran escala. Mientras muchas plataformas modernas delegan ingenuamente reglas críticas a lenguajes interpretados en backend, este proyecto blinda sus reglas de negocio de manera monolítica usando la \textbf{Integridad Referencial y Transaccional} a nivel de Hardware y Datos. Al implementar Particionamiento Físico, Triggers Bloqueantes, Transacciones ACID y Vistas Analíticas, el E-Commerce entrega un rendimiento sumamente robusto y tolerante a fallos lógicos externos, demostrando habilidades a nivel de Diseño Avanzado.

\section{Repositorio de Código (GitHub)}
El código fuente completo del proyecto (incluyendo los esquemas DDL de la base de datos, procedimientos almacenados, triggers, vistas, controladores en Python/Flask y plantillas del frontend) se encuentra alojado y puede ser consultado en la siguiente dirección:

\begin{center}
    \colorbox{backcolour}{
        \parbox{0.95\textwidth}{
            \centering \vspace{0.3cm}
            \textbf{Liga del Repositorio de Git:} \\
            \url{https://github.com/FranciscoSilva28/Proyecto-FInal-Bases-de-Datos-Avanzadas-CRUD-BDA-Steam.git}
            \vspace{0.3cm}
        }
    }
\end{center}

\end{document}
